% Se compila con XeLaTeX; hacen falta varias pasadas para que se resuelvan las referencias.
\documentclass[11pt,a4paper]{article}
\usepackage{fontspec}
\IfFontExistsTF{texgyrepagella-regular.otf}
 {\setmainfont{texgyrepagella}[Extension=.otf,UprightFont=*-regular,BoldFont=*-bold,
   ItalicFont=*-italic,BoldItalicFont=*-bolditalic]}
 {\IfFontExistsTF{Palatino Linotype}{\setmainfont{Palatino Linotype}}{\setmainfont{Latin Modern Roman}}}
\linespread{1.05}
% Con babel-spanish instalado (MiKTeX, TeX Live completo) se usa la opción clásica; si falta,
% se carga el locale español de babel sin patrones de guionado.
\IfFileExists{spanish.ldf}
 {\usepackage[spanish,es-noshorthands,es-nodecimaldot,es-tabla]{babel}}
 {\usepackage{babel}\babelprovide[import,main,hyphenrules=nohyphenation]{spanish}}
\usepackage{amsmath,amssymb,amsthm,mathtools}
\usepackage{unicode-math}
\IfFontExistsTF{texgyrepagella-math.otf}{\setmathfont{texgyrepagella-math.otf}}{\setmathfont{Latin Modern Math}}
\usepackage[margin=26mm]{geometry}
\IfFileExists{microtype.sty}{\usepackage{microtype}}{}
\IfFileExists{titlesec.sty}{\usepackage[nobottomtitles*]{titlesec}
  \renewcommand{\bottomtitlespace}{.12\textheight}}{}
\usepackage{booktabs,array}
\usepackage[shortlabels]{enumitem}
\usepackage{tikz}
\usetikzlibrary{arrows.meta,positioning}
\usepackage{xcolor}
\usepackage[colorlinks=true,linkcolor={blue!45!black},citecolor={green!35!black},urlcolor={blue!55!black}]{hyperref}

\AtBeginDocument{\renewcommand{\tablename}{Tabla}}
\setlength{\parskip}{3pt}
\setlength{\emergencystretch}{2em}
\allowdisplaybreaks[2]
\numberwithin{equation}{section}

\theoremstyle{plain}
\newtheorem{teorema}{Teorema}[section]
\newtheorem{proposicion}[teorema]{Proposición}
\newtheorem{lema}[teorema]{Lema}
\newtheorem{corolario}[teorema]{Corolario}
\theoremstyle{definition}
\newtheorem{definicion}[teorema]{Definición}
\theoremstyle{remark}
\newtheorem{observacion}[teorema]{Observación}

% Notación
\newcommand{\ind}{\mathbf{1}}
\DeclareMathOperator{\Lip}{Lip}
\newcommand{\cA}{{\mathcal{A}}}
\newcommand{\cH}{{\mathcal{H}}}
\newcommand{\cE}{{\mathcal{E}}}
\newcommand{\cG}{{\mathcal{G}}}
\newcommand{\sopt}{s^{*}}
